% ESTRUTURA DO LIVRO
% ==================
% Público: equipe de diagramação. Este módulo controla páginas institucionais
% e a abertura visual dos capítulos. Autores recebem apenas artigo_modelo/.

% Use para ficha técnica, agradecimentos, prefácio e autores enquanto o
% conteúdo final não foi entregue. Não elimine essas chamadas de livro.tex: a
% ordem é exigida pelas diretivas; troque o marcador pelo conteúdo definitivo.
\newcommand{\folhareservada}[1]{%
  \clearpage
  \thispagestyle{empty}%
  \vspace*{\fill}
  \begin{center}\textbf{#1}\\[1cm]\textit{[Conteúdo editorial pendente]}\end{center}
  \vspace*{\fill}
  \clearpage
}

% Ficha catalográfica demonstrativa. Substitua este modelo fictício por uma
% ficha validada antes da publicação e preencha os responsáveis editoriais.
\newcommand{\fichacatalografica}{%
  \clearpage
  \thispagestyle{empty}%
  \vspace*{\fill}
  \begin{center}
    \fbox{%
      \begin{minipage}{0.86\textwidth}
        \small
        \textbf{MODELO FICTÍCIO — FICHA CATALOGRÁFICA}\par
        \vspace{0.4cm}
        \TituloLivro. \AutoresLivro. -- \LocalLivro, \the\year.\par
        \vspace{0.3cm}
        [Número de páginas] p. : il.\par
        \vspace{0.3cm}
        1. Tecnologia. 2. [Assunto]. I. Título.\par
        \vspace{0.3cm}
        CDD: [classificação] \hfill CDU: [classificação]\par
        \vspace{0.5cm}
        \hrule
        \vspace{0.3cm}
        \textbf{Equipes editoriais}\par
        Revisão: [nome(s) dos membros]\par
        Conselho editorial: [nome(s) dos membros]\par
        Diagramação: [nome(s) dos membros]\par
        Modelos: [nome(s) dos responsáveis]
      \end{minipage}%
    }
  \end{center}
  \vspace*{\fill}
  \clearpage
}

% Seção final de autores, inspirada no padrão do Livro - Sata: nome em destaque
% e currículo logo abaixo. Os textos são preenchidos em autores.tex, conforme
% cada autor enviar seu currículo.
\newcommand{\autoreslivro}{%
  \cleardoublepage
  \definircabecalholivro
  \chapter*{AUTORES}
  \thispagestyle{fancy}%
  \addcontentsline{toc}{chapter}{Autores}%
  \markboth{Autores}{Autores}%
  % AUTORES — seção final do livro
% =============================
% Cada autor deve enviar seu próprio currículo para a equipe editorial.
% Substitua o exemplo abaixo por um bloco \perfilautor{Nome}{Currículo} para
% cada autor, seguindo o padrão do Livro - Sata.

\perfilautor{Autor de Exemplo}{%
  Lorem ipsum dolor sit amet, consectetur adipiscing elit. Graduado em
  Tecnologia da Informação pela Instituição de Ensino de Exemplo, possui
  especialização em sistemas computacionais e experiência em projetos de
  pesquisa, extensão e desenvolvimento de soluções digitais. Atua nas áreas
  de análise de sistemas, educação tecnológica e inovação, participando de
  atividades acadêmicas, profissionais e de produção científica. Este texto é
  apenas demonstrativo e deve ser substituído pelo currículo do autor.
}

% Exemplo de preenchimento:
% \perfilautor{Nome completo do autor}{%
%   Titulação, instituição e breve descrição da atuação acadêmica e profissional.
% }
%
  \clearpage
}

\newcommand{\perfilautor}[2]{%
  \par\medskip
  \noindent{\large\bfseries\MakeUppercase{#1}}\par
  \smallskip
  \noindent #2\par
  \medskip
}

% A contracapa deve usar a marca IFPI. A equipe de diagramação pode substituir
% este layout inteiro, desde que preserve a marca e a posição no livro.tex.
\newcommand{\contracapareservada}{%
  \clearpage
  \thispagestyle{empty}%
  \centering
  \marcaifpi
  \vfill
  \textit{[Conteúdo editorial da contracapa pendente]}
  \vfill
  \clearpage
}

% Capa e folha de rosto usam os valores de identidade.tex. Se houver arte final
% fornecida pela diagramação, edite estas macros; livro.tex não precisa mudar.
\newcommand{\capalivro}{%
  \begin{titlepage}
    \centering
    \marcaifpi\\[2cm]
    \vspace*{5cm}
    {\fontsize{24}{28}\selectfont\color{sommagreen}\textbf{\TituloLivro}}\\[1cm]
    {\Large \SubtituloLivro}\\[3cm]
    {\large Autores: \AutoresLivro}
    \vfill
    {\large \LocalLivro\\ \the\year}
  \end{titlepage}
}
\newcommand{\folhaderosto}{%
  \begin{titlepage}
    \centering
    \marcaifpi\\[2cm]
    \vspace*{4cm}
    {\fontsize{20}{24}\selectfont\color{sommagreen}\textbf{\TituloLivro}}\\[1cm]
    {\Large \SubtituloLivro}\\[2cm]
    {\large Autores: \AutoresLivro}\\
    \vfill
    {\large \LocalLivro\\ \the\year}
  \end{titlepage}
}

% Abertura simples: identificação do capítulo, título, autoria, resumo e palavras-chave.
\newcommand{\capitulocapa}[6]{%
  \cleardoublepage
  \refstepcounter{chapter}%
  \setcounter{section}{0}%
  \setcounter{table}{0}%
  \setcounter{figure}{0}%
  \registracapitulonosumario{\thechapter}{#1}{#5}%
  \definircabecalholivro
  \thispagestyle{fancy}%
  \begin{center}
    {\large\bfseries Capítulo \thechapter}\\[0.8cm]
    {\fontsize{14}{16}\selectfont\color{sommagreen}\textbf{#1}\\ #2}\\[0.8cm]
    {\fontsize{12}{14}\selectfont\textbf{#3}}\\[0.4cm]
    \vspace{0.2cm}
    \begin{minipage}{0.88\textwidth}
      \small\textbf{Resumo.} #4\par\vspace{0.5cm}
      \textbf{Palavras-chave:} #6
    \end{minipage}
  \end{center}
  \clearpage
}
